\section{Implementierung und Ergebnisse}

Mittels all dieser Kommandos und Umgebungen ist es uns nun möglich, eine
vollständige Hausarbeit zu schreiben. Dafür müsst ihr das \LaTeX~Projekt 
lediglich herunterladen und könnt quasi sofort loslegen. Wie bereits erwähnt,
liegt in diesem Projekt der Inhalt dieser Anleitung. Es ist euch also ohne 
weiteres möglich, den gesamten Inhalt auch nochmal im Quelltext nachzulesen.\\
In der Infrastruktur der Informatik kann das Projekt ohne Probleme gebaut 
werden und es sind keine weiteren Schritte notwendig. Auch mittels Overleaf
sollte es ohne Probleme möglich sein, das Projekt zu importieren und zu bauen. 
Lediglich lokal auf eurem Rechner müsst ihr einige Programme installieren, zum 
einen braucht ihr \LaTeX~für euer Betriebssystem, zum anderen wird eine Form 
von Editor benötigt. Auf macOS und Linux steht euch der Vim zur Verfügung, auf
Windows empfehle ich die Installation der WSL, da \LaTeX~unter Windows nicht 
unbedingt angenehm zu nutzen ist und mein Projekt auch nicht ohne weiteres unter
Windows gebaut werden kann. Zusätzlich wird noch das Programm \texttt{pygments} 
benötigt.\\
Steht die passende Umgebung zur Verfügung, könnt ihr mein Projekt entweder über 
den Browser oder noch einfacher über die Kommandozeile herunterladen. Die passenden
Befehle und URL findet ihr auf meiner 
\href{https://informatik.hs-bremerhaven.de/flignitz/latex.html}{Webseite}.\\
Um das gesamte Dokument zu bauen (kompilieren) reicht es, das hinterlegte Skript
auszuführen. Grundsätzlich läuft das Kompilieren mehrfach durch, dies ist unter 
\LaTeX~üblich. Leider bietet der Compiler nicht unbedingt sehr aussagekräftige 
Fehlermeldungen, es gilt die Ausgabe genau zu lesen. Ebenfalls kann der Compiler auch
mit fehlerhaftem Code oftmals ein PDF produzieren, dies enthält dann jedoch nicht 
zwingend den gewünschten Inhalt. Ihr solltet entsprechend darauf achten, ob der 
Prozess tatsächlich sauber durchläuft. Dies wird auch am Ende des Kompilierens 
ausgegeben, der Fehlercode \texttt{1} ist dabei das eindeutigste Indiz für einen
Fehler.
